\documentclass[12pt]{article}
\usepackage[T1]{fontenc}
\usepackage[utf8]{inputenc}
\usepackage{lmodern}
\usepackage{textcomp}
\usepackage{enumitem}
\usepackage{geometry}
\geometry{margin=1in}
\begin{document}
\begin{center}
Answer Key - Computer Science - Undergraduate Level Assessment 14 \
\small (Answers are ordered exactly as in the question paper.)
\end{center}

\section*{Multiple Choice}
\begin{enumerate}[label=\arabic*.]
\item \textbf{Answer:} C
\\ \textit{Options:} A) Bubble sort; B) Insertion sort; C) Merge sort; D) Quicksort
\\ \textit{Note:} Merge sort consistently divides the input array in half and recursively sorts each half, leading to an average-case and worst-case time complexity of O(n log n) due to the logarithmic number of divisions and linear time required to merge the sorted halves.
\item \textbf{Answer:} D
\\ \textit{Options:} A) O(log N); B) O(N log N); C) O(log N) + 0(1); D) O((log $N)^2$)
\\ \textit{Note:} The recurrence relation f(2 N + 1) = f(2 N) = f(N) + log N indicates that each recursive call adds a logarithmic term, and since the depth of the recursion grows logarithmically with N, the overall complexity accumulates to O((log $N)^2$).
\item \textbf{Answer:} B
\\ \textit{Options:} A) 3; B) 4; C) 2; D) 8
\\ \textit{Note:} In Python 3, the operator "\textgreater{}\textgreater{}" represents a bitwise right shift, so shifting the binary representation of 8 (which is 1000 in binary) one position to the right results in 4 (which is 0100 in binary).
\item \textbf{Answer:} D
\\ \textit{Options:} A) To translate Web addresses to host names; B) To determine the IP address of a given host name; C) To determine the hardware address of a given host name; D) To determine the hardware address of a given IP address
\\ \textit{Note:} The Address Resolution Protocol (ARP) is designed to map an IP address to its corresponding hardware address (MAC address) within a local network, enabling effective communication between devices.
\item \textbf{Answer:} A
\\ \textit{Options:} A) Replace the page whose next reference will be the longest time in the future.; B) Replace the page whose next reference will be the shortest time in the future.; C) Replace the page whose most recent reference was the shortest time in the past.; D) Replace the page whose most recent reference was the longest time in the past.
\\ \textit{Note:} Replacing the page whose next reference will be the longest time in the future minimizes page faults by ensuring that the most immediately useful pages remain in memory, thereby reducing the likelihood of needing to load a page that will be referenced sooner.
\end{enumerate}

\section*{True / False}
\begin{enumerate}[label=\arabic*.]
\setcounter{enumi}{5}
\item \textbf{Answer:} True
\\ \textit{Options:} A) True; B) False
\\ \textit{Note:} O(1) represents constant time complexity, meaning that the execution time of an algorithm remains constant regardless of the size of the input data, thus making it the smallest asymptotic notation in terms of growth rate.
\item \textbf{Answer:} True
\\ \textit{Options:} A) True; B) False
\\ \textit{Note:} Even after extensive testing, a program may still contain undetected bugs because testing can only cover a finite number of scenarios and may not account for every possible input or interaction within the software.
\item \textbf{Answer:} True
\\ \textit{Options:} A) True; B) False
\\ \textit{Note:} The statement is true because if a \textless{} c is true, then regardless of the values of b and the condition !(a == c), the entire expression a \textless{} c || a \textless{} b \&\& !(a == c) will evaluate to true due to the logical OR operator, which requires only one true operand to yield true.
\item \textbf{Answer:} True
\\ \textit{Options:} A) True; B) False
\\ \textit{Note:} In Python 3, a set automatically removes duplicate elements from a collection, thereby ensuring that only unique elements are retained when converting a list to a set.
\item \textbf{Answer:} True
\\ \textit{Options:} A) True; B) False
\\ \textit{Note:} The statement is true because in binary notation, 0.1 represents 1/2, which is exactly equal to the decimal number 0.5.
\end{enumerate}

\section*{Long Form Response}
\begin{enumerate}[label=\arabic*.]
\setcounter{enumi}{10}
\item \textbf{Answer:} def get\_odd\_occurence(arr, arr\_size): | return arr[i] | return -1
\\ \textit{Note:} correct because it suggests a function that iterates through the array to identify and return the number that appears an odd number of times, although it lacks the necessary implementation details and loop structure to achieve this.
\item \textbf{Answer:} def sort\_sublists(input\_list): | return result
\\ \textit{Note:} correct because it outlines the structure of a Python function that utilizes a lambda function to sort each sublist of strings within a list of lists, aligning with the question's requirement for sorting functionality.
\end{enumerate}

\vfill
\noindent\textit{End of Answer Key}
\end{document}